
\documentclass[a4paper,11pt]{article}
\usepackage[a4paper, margin=8em]{geometry}

% usa i pacchetti per la scrittura in italiano
\usepackage[french,italian]{babel}
\usepackage[T1]{fontenc}
\usepackage[utf8]{inputenc}
\frenchspacing 

% usa i pacchetti per la formattazione matematica
\usepackage{amsmath, amssymb, amsthm, amsfonts}

% usa altri pacchetti
\usepackage{gensymb}
\usepackage{hyperref}
\usepackage{standalone}

% imposta il titolo
\title{Appunti Comunicazioni Numeriche}
\author{Luca Seggiani}
\date{2026}

% imposta lo stile
% usa helvetica
\usepackage[scaled]{helvet}
% usa palatino
\usepackage{palatino}
% usa un font monospazio guardabile
\usepackage{lmodern}

\renewcommand{\rmdefault}{ppl}
\renewcommand{\sfdefault}{phv}
\renewcommand{\ttdefault}{lmtt}

% disponi il titolo
\makeatletter
\renewcommand{\maketitle} {
	\begin{center} 
		\begin{minipage}[t]{.8\textwidth}
			\textsf{\huge\bfseries \@title} 
		\end{minipage}%
		\begin{minipage}[t]{.2\textwidth}
			\raggedleft \vspace{-1.65em}
			\textsf{\small \@author} \vfill
			\textsf{\small \@date}
		\end{minipage}
		\par
	\end{center}

	\thispagestyle{empty}
	\pagestyle{fancy}
}
\makeatother

% disponi teoremi
\usepackage{tcolorbox}
\newtcolorbox[auto counter, number within=section]{theorem}[2][]{%
	colback=blue!10, 
	colframe=blue!40!black, 
	sharp corners=northwest,
	fonttitle=\sffamily\bfseries, 
	title=~\thetcbcounter: #2, 
	#1
}

% disponi definizioni
\newtcolorbox[auto counter, number within=section]{definition}[2][]{%
	colback=red!10,
	colframe=red!40!black,
	sharp corners=northwest,
	fonttitle=\sffamily\bfseries,
	title=~\thetcbcounter: #2,
	#1
}

% U.D.M
\newcommand{\amp}{\ensuremath{\mathrm{A}}}
\newcommand{\volt}{\ensuremath{\mathrm{V}}}
\newcommand{\meter}{\ensuremath{\mathrm{m}}}
\newcommand{\second}{\ensuremath{\mathrm{s}}}
\newcommand{\farad}{\ensuremath{\mathrm{F}}}
\newcommand{\henry}{\ensuremath{\mathrm{H}}}
\newcommand{\siemens}{\ensuremath{\mathrm{S}}}

% circuiti
\usepackage{circuitikz}
\usetikzlibrary{babel}

% disegni
\usepackage{pgfplots}
\pgfplotsset{width=10cm,compat=1.9}

% disponi codice
\usepackage{listings}
\usepackage[table]{xcolor}

\lstdefinestyle{codestyle}{
		backgroundcolor=\color{black!5}, 
		commentstyle=\color{codegreen},
		keywordstyle=\bfseries\color{magenta},
		numberstyle=\sffamily\tiny\color{black!60},
		stringstyle=\color{green!50!black},
		basicstyle=\ttfamily\footnotesize,
		breakatwhitespace=false,         
		breaklines=true,                 
		captionpos=b,                    
		keepspaces=true,                 
		numbers=left,                    
		numbersep=5pt,                  
		showspaces=false,                
		showstringspaces=false,
		showtabs=false,                  
		tabsize=2
}

\lstdefinestyle{shellstyle}{
		backgroundcolor=\color{black!5}, 
		basicstyle=\ttfamily\footnotesize\color{black}, 
		commentstyle=\color{black}, 
		keywordstyle=\color{black},
		numberstyle=\color{black!5},
		stringstyle=\color{black}, 
		showspaces=false,
		showstringspaces=false, 
		showtabs=false, 
		tabsize=2, 
		numbers=none, 
		breaklines=true
}

\lstdefinelanguage{javascript}{
	keywords={typeof, new, true, false, catch, function, return, null, catch, switch, var, if, in, while, do, else, case, break},
	keywordstyle=\color{blue}\bfseries,
	ndkeywords={class, export, boolean, throw, implements, import, this},
	ndkeywordstyle=\color{darkgray}\bfseries,
	identifierstyle=\color{black},
	sensitive=false,
	comment=[l]{//},
	morecomment=[s]{/*}{*/},
	commentstyle=\color{purple}\ttfamily,
	stringstyle=\color{red}\ttfamily,
	morestring=[b]',
	morestring=[b]"
}

% disponi sezioni
\usepackage{titlesec}

\titleformat{\section}
	{\sffamily\Large\bfseries} 
	{\thesection}{1em}{} 
\titleformat{\subsection}
	{\sffamily\large\bfseries}   
	{\thesubsection}{1em}{} 
\titleformat{\subsubsection}
	{\sffamily\normalsize\bfseries} 
	{\thesubsubsection}{1em}{}

% disponi alberi
\usepackage{forest}

\forestset{
	rectstyle/.style={
		for tree={rectangle,draw,font=\large\sffamily}
	},
	roundstyle/.style={
		for tree={circle,draw,font=\large}
	}
}

% disponi algoritmi
\usepackage{algorithm}
\usepackage{algorithmic}
\makeatletter
\renewcommand{\ALG@name}{Algoritmo}
\makeatother

% disponi numeri di pagina
\usepackage{fancyhdr}
\fancyhf{} 
\fancyfoot[L]{\sffamily{\thepage}}

\makeatletter
\fancyhead[L]{\raisebox{1ex}[0pt][0pt]{\sffamily{\@title \ \@date}}} 
\fancyhead[R]{\raisebox{1ex}[0pt][0pt]{\sffamily{\@author}}}
\makeatother

\begin{document}
% sezione (data)
\section{Lezione del 17-04-26}

% stili pagina
\thispagestyle{empty}
\pagestyle{fancy}

% testo
\subsection{Indici caratteristici di distribuzioni note}
Nella sezione 13.1 avevamo visto gli \textit{indici caratteristici} di una distribuzione $X$. Calcoliamoli adesso per distribuzioni note.

\begin{itemize}
\item \textbf{\textsf{Distribuzione gaussiana}} \\
Per la distribuzione gaussiana prendiamo la definizione standard:
$$
f_z(z) = \frac{1}{\sqrt{2 \pi}} e^{-\frac{z^2}{2}}
$$
dove si era posto:
$$
\begin{cases}
\eta_z = E[Z] = 0 \\
\sigma_z^2 = E[(z - \eta_z)^2] = E(z^2) = \int_{-\infty}^{\infty} z^2 \frac{1}{\sqrt{2 \pi}} e^{-\frac{z^2}{2}} \, dz = 1
\end{cases}
$$
Assumendo $\eta_z$ come il valor medio, e $\sigma_z^2$ come la varianza.
Questo si dimostra vero dal fatto che la $f_z(z)$ è pari, per il valor medio, e dal calcolo dell'integrale (che non riportiamo ma ricordiamo si fa passando al caso polare tridimensionale) per la varianza.

Possiamo quindi chiederci come si calcolano valor medio e varianza di una variabile aleatoria gaussiana arbitraria. Ricordiamo innanzitutto la trasformazione:
$$
X = \sigma_x Z + \eta_x
$$
che ci porta dalla definizione standard alla distribuzione gaussiana con $\eta_x$, $\sigma_x$ arbitrari. 
A questo punto possiamo applicare il teorema fondamentale delle trasformazioni di variabili continue (13.1) per dire:
$$
f_x(x) = \sum_{i = 1}^k \frac{f_z(z_i)}{|g'(z_i)|} \Bigg|_{z_i = g^{-1}(x)}
= \frac{f_z(z_1)}{|g'(z_i)|} \Bigg|_{z_i = \frac{x - \eta_x}{\sigma_x}}
= \frac{1}{\sigma_x} \cdot \frac{1}{\sqrt{2 \pi}} e^{-\left( \frac{x - \eta_x}{\sigma_x} \right)^2 \cdot \frac{1}{2}}
$$
in quanto $k = 1$ visto che si ha una sola inversa per ogni ordinata di una retta.
Notiamo quindi che si ottiene la forma che ci aspettiamo per una variabile aleatoria in distribuzione gaussiana con valor medio $\eta_x$ e varianza $\sigma_x^2$. Siamo allora giustificati ad usare il teorema dell'aspettazione:
$$
\begin{cases}
E[X] = \sigma_x E(z) + \eta_x = \eta_x \\
E[(X - \eta_x)^2] = E[(\sigma_X Z + \eta_x - \eta_x)^2] = E[(\sigma_x Z)^2] 
= \sigma_x^2 \cdot E[Z^2] = \sigma_x^2
\end{cases}
$$
e verificare che effettivamente valor medio e varianza sono quelli effettivi calcolati dalle definizioni date sugli indici.
\end{itemize}

\subsubsection{Assenza di memoria delle variabili aleatorie esponenziali}
Vediamo una proprietà particolare della variabili aleatorie esponenziali. Queste erano definite da densità e distribuzione di probabilità:
$$
f_X(x) = \frac{1}{\lambda} e^{-\frac{x}{\lambda}} u(x), \quad
F_X(x) = \left( 1 - e^{-\frac{x}{\lambda}} \right) u(x)
$$

Chiediamoci se la probabilità oltre un tempo $\Delta t$, a partire da un tempo $t_0$, sapendo di essere arrivati fino al tempo $t_0$, dipende da $t_0$ stesso. In simboli, vogliamo calcolare:
$$
P\{ X > t_0 + \Delta t | X > t_0 \} = \frac{ P\{ X > t_0 + \Delta t \} \cap P\{ X > t_0 \} }{ P\{ X > t_0 \} } = \frac{ P\{ X > t_0 + \Delta t \} }{ P\{ X > t_0 \} }
$$
$$
= \frac{1 - F_X(t_0 + \Delta t)}{1 - F_X(t_0)} 
= \frac{e^{-\frac{t_0 + \Delta t}{\lambda}}}{e^{-\frac{t_0}{\lambda}}}
= \frac{e^{-\frac{t_0}{\lambda}}e^{-\frac{\Delta t}{\lambda}}}{e^{-\frac{t_0}{\lambda}}} = e^{-\frac{\Delta t}{\lambda}} = P\{X > \Delta t\}
$$
Ebbene abbiamo scoperto che la variabile non dipende da $t_0$, ma anzi è interamente dipendente da $\Delta t$, qualsiasi sia il tempo $t_0$ correntemente trascorso. Chiamiamo questa proprietà \textit{assenza di memoria} della variabile aleatoria esponenziale.

\subsection{Teorema di Tchebycheff}
Abbiamo detto che la varianza $\sigma_x^2$ di una variabile aleatoria è una misura della dispersione per una variabile aleatoria, e affinché una variabile aleatoria sia valida sarà $\sigma_x^2 < \infty$.

Se escludiamo un intervallo di $2 k$ deviazioni standard attorno al valor medio $\eta_x$, cioè prendiamo l'intervallo di probabilità:
$$
P\left\{ |X - \eta_x| \geq k \sigma_x \right\}
$$
possiamo dimostrare che questa probabilità ha un limite superiore attraverso il \textbf{teorema di Tchebycheff}:
\begin{theorem}{Teorema di Tchebycheff}
Data una qualsiasi variabile aleatoria $X$ con varianza finita, per ogni $k > 1$ vale:
$$
P\left\{ |X - \eta_x| \geq k \sigma_x \right\} \leq \frac{1}{k^2}
$$
\end{theorem}

Questo teorema ci dice che, per le proprietà della funzione probabilità (dagli assiomi di Kolmogorov), possiamo dire:
$$
P\left\{ |X - \eta_x| < k \sigma_x \right\} = 1 - P\left\{ |X - \eta_x| \geq k \sigma_x \right\} \geq 1 - \frac{1}{k^2}
$$
cioè abbiamo un limite inferiore per un qualsiasi intervallo di $k$ deviazioni standard attorno al valor medio di una variabile aleatoria $X$.

Una definizione alternativa di questo teorema ci dice che, per un qualsiasi intervallo di ampiezza $2 \epsilon$ attorno al valor medio, possiamo determinare il valor minimo della probabilità che la variabile aleatoria $X$ cada in tale intervallo come:
$$
P\left\{ |X - \eta_x| < \epsilon \right\} \geq 1 - \frac{\sigma_x^2}{\epsilon^2}
$$
cioè riscriviamo l'intervallo in $\epsilon$ come un numero di deviazioni standard:
$$
\epsilon = k \sigma_x
$$
e applichiamo la proprietà appena trovata come corollario del teorema di Tchebycheff.

\subsection{Sistemi di più variabili aleatorie}
Possiamo avere sistemi di più variabili aleatorie. Prendiamo ad esempio 2 variabili $X$ ed $Y$. 

\par\smallskip

\noindent
\textbf{\textsf{Probabilità congiunta}}

\noindent
Poiché $X$ e $Y$ sono variabili aleatorie, gli insiemi $\{X \leq x\}$ e $\{Y \leq y\}$ costituiscono eventi, e lo costituisce anche l'intersezione:
$$
\{X \leq x\} \cap \{Y \leq y\}
$$
di cui calcoliamo la probabilità come funzione \textbf{distribuzione di probabilità congiunta}.
\begin{definition}{Distribuzione di probabilità congiunta}
Per due variabili aleatorie $X$ e $Y$, si definisce la distribuzione di probabilità congiunta:
$$
F_{XY}(x) = P\{X \leq x, Y \leq y\}
$$
\end{definition}
La probabilità congiunta è definita come sopra e consiste nella probabilità che le 2 variabili cadano, appunto \textit{congiuntamente}, all'interno dell'intervallo specificato.
Riguardo alla probabilità congiunta valgono le proprietà:
\begin{enumerate}
	\item $0 \leq F_{XY}(x, y) \leq 1$ in quanto si parla comunque di probabilità;
	\item $F_{XY}$ è funzione monotona non descrescente di \textit{ciascuna} delle 2 variabili;
	\item Valgono i limiti, sulle 2 variabili:
\[
\begin{cases}
	F_{XY}(-\infty, -\infty) = F_{XY}(-\infty, 0) = F_{XY}(0, -\infty) = 0 \\
	F_{XY}(\infty, \infty) = 1
\end{cases}
\]
\end{enumerate}

\par\smallskip

\noindent
\textbf{\textsf{Regole marginali}}

\noindent
Le funzioni di distribuzione di ciascuna delle due variabili aleatorie $X$ e $Y$ si ottengono dalle \textbf{regole marginali}:
$$
F_X(x) = F_{XY}(x, \infty), \quad F_Y(y) = F_{XY}(\infty, y)
$$
Questo si ricava, come sopra, dal fatto che:
$$
\{X \leq x, Y \leq \infty\} = \{X \leq x\}, \quad \{X \leq \infty, Y \leq y\} = \{Y \leq y\}
$$

\par\smallskip

\noindent
\textbf{\textsf{Funzione densità di probabilità congiunta}}

\noindent
Dalla la funzione distribuzione di probabilità congiunta, come per il caso monodimensionale, si puà calcolare una funzione \textbf{densità di probabilità congDistribuzioneiunta}: 
\begin{definition}{Densità di probabilità congiunta}
Per due variabili aleatorie $X$ e $Y$, si definisce la densità di probabilità congiunta:
$$
f_{XY} = \frac{\partial^2 F_{XY}(x, y)}{\partial x \partial y}
$$
\end{definition}
Questo, chiaramente, ammesso che la distribuzione sia derivabile fino al secondo grado. Varrà quindi:
$$
f_{XY} (x, y) \, dx dy = P \{ x < X < x + dx, y < Y + y + dy \}
$$
e in particolare:
$$
P[(X, y) \in D] = \int_D f_{XY} (x, y) \, dx dy
$$

Questa sarà quindi una funzione di 2 variabili $(x, y)$, corrispondenti al valore che vogliamo valutare delle variabili aleatorie $X$ e $Y$.

Anche dalla densità di probabilità congiunta si possono ricavare le \textit{funzioni marginali}, stavolta di densità:
$$
f_X(x) = \int_{-\infty}^\infty f_{XY} (x, y) \, dx, \quad
f_X(x) = \int_{-\infty}^\infty f_{XY} (x, y) \, dy
$$
cioè semplicemente si ottengono le funzioni densità di probabilità in $X$ o in $Y$ rispettivamente integrando su $dx$ o su $dy$. 

\subsubsection{Sistemi di variabili aleatorie indipendenti}
Ricordiamo la definizione di indipendenza fra eventi aleatori da 5.1.3. Questa diceva che 2 eventi sono detti indipendenti se:
$$
P(a \cap b) = P(a) P(b)
$$
Ne possiamo dare una definizione anche per le probabilità congiunto prendendo:
\begin{definition}{Sistemi di variabili aleatorie indipendenti}
In un sistema di 2 variabili aleatorie, $X$ e $Y$, si dice	che $X$ e $Y$ sono indipendenti se vale rispetto alle distribuzione di probabilità congiunta e alle distribuzioni di probabilità di $X$ e $Y$:
$$
F_{XY}(x, y) = P\{X \leq x, Y \leq y\} = P\{X \leq x\}\{Y \leq y\} = F_X(x) F_Y(y)
$$
\end{definition}
Da questa definizione segue che anche la densità di probabilità congiunta si fattorizza nel prodotto delle due densità di probabilità marginali:
$$
f_{XY}(x, y) = \frac{\partial^2 F_{XY}(x, y)}{\partial x \partial y} = \frac{\partial F_X (x)}{\partial x} \cdot \frac{\partial F_Y (y)}{\partial y} = f_X(x) f_Y(y)
$$
Osserviamo quindi che nel caso di variabili aleatorie indipendenti le densità di probabilità marginali sono sufficienti per descrivere statisticamente l'intero sistema.

\end{document}
